%========================
% Document class and theme
%========================
\documentclass[8pt]{beamer}
\usetheme[progressbar=frametitle]{metropolis}
\setbeamersize{text margin left=10mm, text margin right=10mm}
\usepackage{appendixnumberbeamer} % appendix slide numbering
\setbeamertemplate{theorems}[numbered]

%========================
% Core packages
%========================
\usepackage{amsmath, amsfonts, amssymb, amsthm} % math + theorems
\usepackage{booktabs}        % professional tables
\usepackage{hyperref}        % hyperlinks
\usepackage{xcolor}          % colors
\usepackage{xspace}          % spacing for custom commands

%========================
% Algorithms
%========================
\usepackage{algorithm}
\usepackage{algpseudocode}
\newtheorem{proposition}{Proposition}
\usepackage{bbm}

%========================
% Plots and TikZ
%========================
\usepackage{pgfplots}
\usepgfplotslibrary{dateplot}
\usepackage{tikz}
\usetikzlibrary{positioning}

%========================
% Listings (code)
%========================
\usepackage{listings}
\lstset{
    basicstyle=\ttfamily\small,
    breaklines=true,
    numbers=left,
    numberstyle=\tiny
}

% R style
\lstdefinelanguage{R}{
  morekeywords={TRUE,FALSE},
  deletekeywords={data,frame,length,as,character},
  otherkeywords={0,1,2,3,4,5,6,7,8,9},
  keywordstyle=\color{blue},
  commentstyle=\color{DarkGreen},
  stringstyle=\color{DarkGreen},
  basicstyle=\ttfamily\small
}

% Python style
\lstdefinelanguage{PythonCustom}{
  language=Python,
  keywordstyle=\color{blue},
  commentstyle=\color{gray},
  stringstyle=\color{red},
  basicstyle=\ttfamily\small
}

%========================
% Custom commands
%========================
\newcommand{\themename}{\textbf{\textsc{metropolis}}\xspace}

%========================
% Custom footline
%========================
\setbeamertemplate{footline}
{%
  \leavevmode%
  \hbox{%
  \begin{beamercolorbox}[wd=.35\paperwidth,ht=2.5ex,dp=1.5ex,center]{author in head/foot}%
    \usebeamerfont{author in head/foot}\insertshortauthor
  \end{beamercolorbox}%
  \begin{beamercolorbox}[wd=.3\paperwidth,ht=2.5ex,dp=1ex,center]{title in head/foot}%
    \usebeamerfont{title in head/foot}\insertshorttitle
  \end{beamercolorbox}%
  \begin{beamercolorbox}[wd=.3\paperwidth,ht=2.5ex,dp=1ex,right]{date in head/foot}%
    \usebeamerfont{date in head/foot}\insertframenumber{} / \inserttotalframenumber
  \end{beamercolorbox}}%
  \vskip0pt%
}

%========================
% Beamer tweaks
%========================
\setbeamertemplate{navigation symbols}{} % remove default navigation symbols


%%%%%%%%%%%%%%%%%%%%%%%%%%%%%%%%%%%%%%%%%%%%%%%%%%%%%%%%%%%%%%%%%%%%
%%%%%%%%%%%%%%%%%%%%%%%%%%%%%%%%%%%%%%%%%%%%%%%%%%%%%%%%%%%%%%%%%%%%
% AQUI SE DEFINEN LAS IMAGENES PARA UTILIZAR DESPUES
%\pgfdeclareimage[interpolate=true, height=7cm,width=16cm]{halton-points}{halton-points}
%\pgfdeclareimage[interpolate=true, height=3cm, width =4cm]
%{serie-petroleo-reducido}{serie-petroleo-reducido}
%\pgfdeclareimage[interpolate=true, height=3cm, width =4cm]{rectangle-triangle}{rectangle-triangle}
%\pgfdeclareimage[interpolate=true, height=3cm, width =4cm]{any-angle}{any-angle}
%\pgfdeclareimage[interpolate=true, height=3cm, width =4cm]{Pythagoras}{Pythagoras}


\title{Chapter 5 - Beyond Monte Carlo}
\subtitle{Bootstrap Method for Bias and Standard Deviation Estimation.}
\author{Prof. Alex Alvarez, Ali Raisolsadat}
\institute{School of Mathematical and Computational Sciences \\ University of Prince Edward Island}
\date{} % leave empty or add \today
%\title[Stat 4110]{Stat 4110 Statistical Simulation}
%\subtitle{}
%\author[University of Prince Edward Island]{School of Mathematical and Computational Sciences \\ University of Prince Edward Island}

%========================
% Begin document
%========================
\begin{document}

%-------------------
% Title frame
%-------------------
\maketitle

%-------------------
% Slide 1: Framework
%-------------------
\begin{frame}{Framework}

\textbf{Sample Data:}  
Let 
\[
X = (X_1, X_2, \ldots, X_n)
\]
be independent and identically distributed random variables with
\[
X_i \sim F
\]
where the distribution $F$ is unknown. We observe a realization
\[
x = (x_1, x_2, \ldots, x_n)
\]
We are interested in estimating a parameter
\[
\theta = \theta(F)
\]
Assume that we have an estimator 
\[
\widehat{\theta}_n = T(X)
\]
of $\theta$.

We focus on the following problems:
\begin{itemize}
\item Estimation of the bias of $\widehat{\theta}_n$
\item Estimation of the standard deviation of $\widehat{\theta}_n$
\end{itemize}

\end{frame}

%-------------------
% Slide 2: Formulation
%-------------------
\begin{frame}{Formulation}

In our previous lecture we defined the empirical cumulative distribution function:
\[
F_n(t)=\frac{1}{n}\sum_{i=1}^n 1_{(-\infty,t]}(x_i)
\]

The empirical distribution $F_n$ approximates the true distribution $F$, i.e.,
\[
F_n \approx F
\]

We generate Bootstrap samples
\[
x^{*(j)} = \left(x^{*(j)}_1, x^{*(j)}_2, \ldots, x^{*(j)}_n\right), 
\qquad j=1,2,\ldots,N
\]
where the values $x^{*(j)}_m$ are i.i.d. random variables sampled from $F_n$.

An estimator $\widehat{\theta}_n = T(X)$ satisfies the \textbf{plug-in principle} if
\[
\widehat{\theta}_n(x) = \theta(F_n)
\]
for all $x=(x_1,x_2,...,x_n)$.

\end{frame}

%-------------------
% Slide 3: Bias Estimation
%-------------------
\begin{frame}{Bias Estimation}
We know that $bias\left(\widehat{\theta}_n\right)=E(\widehat{\theta}_n)-\theta=E_F(\widehat{\theta}_n)-\theta(F)$

\vspace{3mm}

As we do not know $F$, but we do know $F_n$, we will approximate $bias\left(\widehat{\theta}_n\right)$ with
\begin{equation*}
E_{F_n}(\widehat{\theta}_n)-\theta(F_n)
\end{equation*}

If $\theta_n$ satisfies the plug-in principle, we can replace $\theta(F_n)$ with $\widehat{\theta}_n(x)$. For the estimations of $E_{F_n}(\widehat{\theta}_n)$ we will use the sample mean of the Bootstrap sample estimates $\widehat{\theta}_n(x^{*(j)})$ with $j=1,2,...N$. Then we will have:
\begin{equation*}
\widehat{bias}\left(\widehat{\theta}_n\right) =  \frac{1}{N} \sum_{j=1}^N \widehat{\theta}_n\left(x^{*(j)}\right) - \widehat{\theta}_n(x)
\end{equation*}
\end{frame}

%-------------------
% Slide 4: Bootstrap Bias Estimation Example
%-------------------
\begin{frame}{Bootstrap Bias Estimation Example}

The following data represent students' grades in a course. The grades are assumed to be independent observations from an unknown distribution $F$. 

We use the sample median to estimate the population median and estimate the bias of this estimator using the Bootstrap method.

\vspace{3mm}

\textbf{Grades:}
\[
x = (84, 63, 92, 89, 80, 47, 68, 66, 56)
\]

\vspace{3mm}

\textbf{Solution:}
By ordering the data in ascending order:
\[
(47, 56, 63, 66, 68, 80, 84, 89, 92)
\]
the sample median is
\[
\widehat{\theta}_9(x) = 68
\]
To estimate the bias, we generate $N = 10000$ Bootstrap samples.

\end{frame}
%-------------------
% Slide 5: Example
%-------------------
%-------------------
% Slide 5: Bootstrap Bias Estimation Example
%-------------------
\begin{frame}[fragile]{Bootstrap Bias Estimation Example}

\textbf{R Code}

\begin{lstlisting}
x <- c(84, 63, 92, 89, 80, 47, 68, 66, 56)
N <- 10000
median_vector <- numeric(N)
theta_hat <- median(x)

for(i in 1:N) {
  Y <- sample(x, size=length(x), replace=TRUE) 
  median_vector[i] <- median(Y)
}
average_medians <- mean(median_vector)
bootstrap_bias <- average_medians - theta_hat
\end{lstlisting}

By running this code once, we obtain
\[
\widehat{bias}(\widehat{\theta}_n) = \frac{1}{N} \sum_{j=1}^N \widehat{\theta}_n(x^{*(j)}) - \widehat{\theta}_n(x) \approx 3.7
\]
\end{frame}

%-------------------
% Slide 6: Estimation of Standard Deviation
%-------------------
\begin{frame}{Bootstrap Estimation of Standard Deviation}

The standard deviation of the estimator $\widehat{\theta}_n$ under the unknown distribution $F$ is
\[
s.d._F(\widehat{\theta}_n)
\]
Since $F$ is unknown, we approximate it using the empirical distribution $F_n$:
\[
s.d._F(\widehat{\theta}_n)
\approx
s.d._{F_n}(\widehat{\theta}_n)
\]
We estimate $s.d._{F_n}(\widehat{\theta}_n)$ using the standard deviation of the Bootstrap estimates
\[
\widehat{\theta}_n(x^{*(1)}),
\widehat{\theta}_n(x^{*(2)}),
\ldots,
\widehat{\theta}_n(x^{*(N)})
\]
Therefore, the Bootstrap estimator of the standard deviation is
\[
\widehat{s.d.}(\widehat{\theta}_n)=\sqrt{\frac{1}{N-1}\sum_{j=1}^N\left(\widehat{\theta}_n(x^{*(j)})-\overline{\widehat{\theta}_n^*}\right)^2}
\]
where
\[
\overline{\widehat{\theta}_n^*}=\frac{1}{N} \sum_{j=1}^N \widehat{\theta}_n(x^{*(j)})
\]

\end{frame}

%-------------------
% Slide 7: Bootstrap Standard Deviation Example
%-------------------
\begin{frame}[fragile]{Bootstrap Standard Deviation Example}

\textbf{R Code}

\begin{lstlisting}
x <- c(84, 63, 92, 89, 80, 47, 68, 66, 56)
N <- 10000
median_vector <- numeric(N)

for (i in 1:N) {
  Y <- sample(x, size=length(x), replace=TRUE) 
  median_vector[i] <- median(Y)
}

bootstrap_sd <- sd(median_vector)
\end{lstlisting}

By running this code once, we obtain

\[
\widehat{s.d.}(\widehat{\theta}_n) \approx 8.38
\]

\end{frame}

%-------------------
% Slide 8: Homework
%-------------------
\begin{frame}{Homework}

Using the data from the previous example, define the hybrid estimator
\[
\widehat{\theta}_n=\frac{1}{2}(\text{sample mean} + \text{sample median})
\]
Use the Bootstrap method, with $N=1000$ samples, to estimate the standard deviation of this estimator.
\end{frame}


\begin{frame}
\begin{algorithm}[H]
\caption{Bootstrap Standard Deviation for Hybrid Estimator}
\begin{algorithmic}[1]

\Require Data $x=(x_1,x_2,\dots,x_n)$, number of bootstrap samples $N$

\For{$j = 1$ to $N$}
\State Generate bootstrap sample
\[
x^{*(j)} = (x^{*(j)}_1, x^{*(j)}_2, \dots, x^{*(j)}_n)
\]
by sampling with replacement from $x$

\State Compute bootstrap estimator
\[ \widehat{\theta}_n^{*(j)} = \frac{1}{2}
\left( \text{mean}(x^{*(j)}) + \text{median}(x^{*(j)}) \right) \]
\EndFor
\State Compute bootstrap mean
\[
\overline{\widehat{\theta}_n^*}
=
\frac{1}{N}
\sum_{j=1}^N
\widehat{\theta}_n^{*(j)}
\]
\State Compute bootstrap standard deviation
\[
\widehat{s.d.}(\widehat{\theta}_n) = \sqrt{ \frac{1}{N-1} \sum_{j=1}^N \left( \widehat{\theta}_n^{*(j)} - \overline{\widehat{\theta}_n^*} \right)^2 }
\]

\State \textbf{Return} $\widehat{s.d.}(\widehat{\theta}_n)$

\end{algorithmic}
\end{algorithm}
\end{frame}

\end{document} 


